\cotefile{57}
\scanpage{242}\mspage{- 16 -}
n$^{\mathrm{o}}$ 227

3. \uline{L'application rationnelle définie par un système linéaire}.

On continue à supposer $V$ complète. Soit $\handwritten{\mathscr L}$ un système linéaire,
$X_{\unsure{0|o}}\handwritten{\in\mathscr L}$, et $L_0$ le module dérivé de $\handwritten{\mathscr L}$ par $X_0$. Soit $f_0,\ldots,f_n$ une
base de ce module. Alors les fonctions $(\handwritten{1},f_{\handwritten{1}}/f_0\unsure{,|;}\ldots,f_n/f_0)$ peuvent
être prises comme coordonnées non homogènes d'un point dans l'espace
projectif, qui représente un point générique d'une variété projective
$V'$. Si $t$ est un élément transcendant sur tout, alors $(tf_0,\ldots,tf_n)$
sont les coordonnées homogènes d'un point générique de $V'$. Elles
définissent une application rationnelle $T:V\handwritten{\longrightarrow}V'$ de $V$ sur $V'$.
Un changement de base dans $L_0$ donne une transformation projective de $V'$.

Si on employait un module $L_{\handwritten{1}}$ dérivé de $\handwritten{\mathscr L}$ par un autre diviseur
$X_{\handwritten{1}}$, on voit que l'application rationnelle $T_{\handwritten{1}}$ obtenue comme ci-dessus
diffère de $T$ par une transformation \deleted{linéaire}\inserted{\handwritten{\unsure{projective}}}. On peut donc dire que $T$
est uniquement déterminée par $\handwritten{\mathscr L}$. De plus, si $\handwritten{\mathscr L}$ et $\handwritten{\mathscr L'}$ sont deux
\margin{\handwritten{\(\big|\)}}systèmes linéaires qui ne diffèrent que par une composante fixe, alors
on voit immédiatement qu'ils donnent lieu à la même application rationnelle
$T$ de $V$, les modules dérivés étant les mêmes. Pour étudier l'application
rationnelle associée à un système linéaire, on peut donc se
borner, sans perte de généralité, aux systèmes linéaires sans composantes
fixes.

On va maintenant donner un critère pour que l'application rationnelle
$T$ soit définie en un point $P$ de $V$.

Un point $P$ de $V$ est dit un \uline{point de base} du système linéaire $\handwritten{\mathscr L}$
si $P\handwritten{\in\unsure{\gamma|r}}\handwritten{|}X\handwritten{|}$ pour tout $X\handwritten{\in\mathscr L}$. Si $Y$ est une composante fixe de $\handwritten{\mathscr L}$, alors
tous les points de $\handwritten{|}Y\handwritten{|}$ sont des points de base. Il peut très bien
arriver qu'un système linéaire n'ait pas de composante fixe, mais ait
% archive sheet number: 241

\scanpage{243}
\[
\operatorname{Pic}(X_n)\overset{\sim}{\longrightarrow}\underline{\operatorname{Pic}}_{X_n/k}(k)
\]
et du th. précédent.

Dans le cas général, Notons \deleted{\unsure{d'abord}} ceci [\deleted{\unsure{si|où}} $k'/k$ est une extension
finie de $k$ \unsure{(G)}, \unsure{qu'on|bien} \unsure{choisira|choisie} de façon

\deleted{\uline{Lemme} \unsure{Sous|Sans} les conditions du corollaire
[car. 0 \unsure{essentielle|considérée}, mais on \unsure{le rappelle|la rappelle}],
le \unsure{conoyau|noyau} de $\operatorname{Pic}(X_{m})\unsure{-|\longrightarrow}\operatorname{Pic}(X_n)$
(pour $m\geq n$) est sans torsion.

On est ramené au cas où $m=n+1$,
où cela \deleted{\unsure{se}} résulte du fait que ce
conoyau est $\subset H^2(X_0,\deleted{\unsure{\mathscr F_n|F_n}}\inserted{F^{(n)}})$, $F^{(n)}_{\deleted{\unsure{n}}}$
un Module \unsure{coh|coh.} convenable sur $X_0$.}

que \deleted{\unsure{les|le}} $X_0\otimes_k k'$ \deleted{\unsure{soit}}satisfasse à la condition
énoncée plus haut] : on a une suite exacte

% slanted
\margin{\uline{\unsure{référence|références|vérifions}}?? \uline{À expliciter}}
\[
0\longrightarrow\operatorname{Pic}(X_n)\longrightarrow\underline{\operatorname{Pic}}_{X_n/k}(k)\longrightarrow\check H^2(k'/k,M_n)
\]
où $M_n$ est le schéma en groupes [comm. lisses
connexes] \deleted{\unsure{tel|tels}} qui décrit les unités
de la $k$-algèbre \deleted{\unsure{$H$}} $H^0(X_n,\mathcal O_{X_n})$,

On note que le système projectif des
$(M_n)$ satisfait M.L. [groupes \deleted{\unsure{alg. connexes}}\inserted{\uline{de type fini} sur $k$ !}]
\unsure{et est|le est} : noyaux \deleted{\unsure{unipotents}} \unsure{finis|lisses|tous}, \uline{vectoriels}\ldots]
% archive sheet number: 242

\scanpage{244}\mspage{- 15 -}
n$^{\mathrm{o}}$ 227

linéaire, le système des hypersurfaces de degré donn\handwritten{\'}{e} sur une
variété projective.

Soit donc $V$ une variété projective, et $R=k[x]=k[x_{\unsure{0|o}},\ldots,x_N]$
son anneau de coordonnées homogènes, comme dans IAG, ch.IV. Soit $I$
\deleted{la fermeture intégrale de $R$}, et $R_n$, $I_n$ les éléments homogènes de
degré $n$ de $R$ et de $I$. Soit $f\handwritten{\in}R_n$. On va associer un diviseur positif
$Z(f)$ à $f$, qui sera appelé l'hypersurface section de $V$ par $f$. Soit
$W^{r-\handwritten{1}}$ un diviseur premier de $V$. Supposons que $W$ ne soit pas contenu
dans l'hyperplan $X_0=0$. Alors $f(x)/x_0^n$ est une fonction $\handwritten{\varphi}$ sur $V$. On
définit la multiplicité de $W$ dans le diviseur $Z(f)$ comme étant
$\operatorname{ord}_W\handwritten{\varphi}$. (Il est clair que $\operatorname{ord}_W\handwritten{\varphi}$ ne dépend pas de l'hyperplan auxiliaire,
et qu'on obtient bien un diviseur positif). Les assertions
suivantes découlent trivialement des définitions :

Tous les cycles $Z(f)$, $f\handwritten{\in}R_n$, sont linéairement équivalent\unsure{:|.} En
fait, si $f$, $g\handwritten{\in}R_n$, on a $(f/g)=Z(f)-Z(g)$.

Si $f_0\handwritten{\in}R_m$, alors l'ensemble des fonctions $f/f_0$, $f\handwritten{\in}R_m$ est
un module $L_0$, et le sup des diviseurs polaires de ces fonctions est
exactement $Z(f_0)$. L'ensemble des diviseurs $Z(f)$ où $f$ parcourt $R_n$ est
un système linéaire, définit par le module $L_0$. On voit aussi qu'il
n'a pas de composantes fixes.

Le système linéaire complet dérivé du système d'hypersurface
de degré $n$ est caractérisé algébriquement par le théorème suivant
(Zariski, th.1 de ES, démonstration facile et élémentaire ) :

\margin{\handwritten{?}}TH\handwritten{\'}{E}OR\handwritten{\`}{E}ME 2. - \uline{Si} $\handwritten{\varphi}$ \uline{est une fonction telle que} $(\handwritten{\varphi})\handwritten{\geq}-Z(f)$ \uline{pour}
$f\handwritten{\in}R_n$, \uline{alors il existe} $g\handwritten{\in}I_n$ \uline{tel que} $\handwritten{\varphi}=g/f$. \uline{Si} $V$ \uline{est normale},
\uline{alors le système d'hypersurfaces de degré} $n$ \uline{est complet pour tout} $n$
\uline{suffisam\inserted{m}ent grand}. (Voir la prop.11 de IAG, ch.IV).
% archive sheet number: 243

\scanpage{245}
% paper-clip impression at upper left
donc pour $n$ grand \inserted{$(\geq n_0)$}, et $m\geq n$,
\[
\check H^{\unsure{*|\times}}(k'/k,M_m)\overset{\sim}{\longrightarrow}\check H^*(k'/k,M_n)
\]

\deleted{Or on a un homomorphisme}

\deleted{D'où facilement} Donc si
$\unsure{\xi|\zeta}_{\unsure{m|n}}\in\underline{\operatorname{Pic}}_{X_m/k}(k)$ est tel que son image
$\xi_n$ dans $\underline{\operatorname{Pic}}_{X_n/k}(k)$ provient de $\underline{\operatorname{Pic}}(X_n)$, alors
il provient lui-$=$ de $\underline{\operatorname{Pic}}(X_m)$. Ceci, et
le fait que $(\underline{\operatorname{Pic}}_{X_n/k}(k))$ satisfait à
M.L. (A.R., montre que il en est
de $=$ de $(\operatorname{Pic}(X_n))_n$, de façon
précise si $n\geq n_0$, \deleted{\unsure{alors il y a|tout élément}}
et si \deleted{\unsure{$m$ est}}\inserted{$n'$} est tel que $\operatorname{Im}(\underline{\operatorname{Pic}}_{X_{\deleted{\unsure{m}}n'}/k}(k)\unsure{-|\to}\underline{\operatorname{Pic}}_{X_n/k}(k)\unsure{)}$
soit \unsure{minimal|minimale|universel}, alors $\operatorname{Im}(\operatorname{Pic}(X_{\deleted{\unsure{m}}n'})\to\operatorname{Pic}(X_n))$
est \unsure{minimal|minimale|universel}. OK.

\diagram{\begin{tikzcd}[ampersand replacement=\&,column sep=small,row sep=small]
0\arrow[r,dash] \& \operatorname{Pic}(X_n)\arrow[d,phantom,"\rotatebox{-45}{$\ni$}"]\arrow[r,dash] \& \underline{\operatorname{Pic}}_{X_n/k}(k)\arrow[d,phantom,"\rotatebox{-45}{$\ni$}"]\arrow[r,dash] \& \deleted{\check H^2(k'/k,M_n)}\\
\& \xi_n\arrow[r,decorate,decoration={snake,amplitude=.35mm,segment length=2mm,post length=1mm}] \& \bar\xi_n \& \\
0\arrow[r,dash] \& \operatorname{Pic}(X_{n'})\arrow[uu]\arrow[r]\arrow[d,phantom,"\rotatebox{-45}{$\ni$}"] \& \underline{\operatorname{Pic}}_{X_{n'}/k}(k)\arrow[uu] \& \\
\& \xi_{n'}\arrow[uu,decorate,decoration={snake,amplitude=.35mm,segment length=2mm,post length=1mm}] \& {} \& \\
0\arrow[r,dash] \& \operatorname{Pic}(X_{n''})\arrow[uu]\arrow[r,dash]\arrow[d,phantom,"\rotatebox{-45}{$\ni$}"] \& \underline{\operatorname{Pic}}_{X_{\unsure{k''|n''}}/k}(k)\arrow[uu]\arrow[d,phantom,"\rotatebox{-45}{$\ni$}"] \& \\
\& \deleted{\unsure{\xi_m}}\lambda_m\arrow[r,decorate,decoration={snake,amplitude=.35mm,segment length=2mm,post length=1mm}] \& \unsure{\eta|\zeta}_m\arrow[uuu,decorate,decoration={snake,amplitude=.35mm,segment length=2mm,post length=1mm},bend right=20] \&
\end{tikzcd}}

\uline{Corollaire} Pour que
$\mathfrak X/\unsure{\mathcal Y|\mathfrak Y|Y}$ soit \uline{projectif},
il f. et s. que pour
tt $n$, $X_n/Y_n$ \inserted{le} soit.
% archive sheet number: 244

\scanpage{246}
Ma chère\\

% The following block is turned through 180 degrees.
\deleted{\uline{Questions} Soit $H\subset G$, $H$ simple sur $k$, contenant
un \deleted{\unsure{tore maxi|tore maximal}} conjugué de $T$, \deleted{Alors} \unsure{géométriquemt|géométriquement|géométrique}.
Soit $S$ un tore dans $H$. Alors $S$ est contenu
dans un tore maximal contenu dans $H$.}
% archive sheet number: 245

\scanpage{248}
\S~1 Un théorème de géométrie analytique.

Rappelons d'abord deux résultats connus

\uline{Lemme} 1.1. \unsure{Soient|Soit} $X$ un espace \deleted{\unsure{para|pour}} topologique, $Y$
une partie fermée de $X$ ayant un syst. \unsure{fondamental|fond.mental}
de voisinages paracompacts, $F$ un faisceau abélien
sur $X$, $F|Y$ le faisceau induit sur $Y$. Alors
l'hom. can. \deleted{\unsure{est}} suivant est un isom :
\[
\underset{U}{\underrightarrow{\lim}}H^*(U,F)\longrightarrow\unsure{-|\vphantom{-}}H^*(Y,F|Y)
\]
[où dans le premier \unsure{mb|mt|membre}, la \uline{lim} est sur
l'ensemble des voisinages de $Y$ \unsure{ds|de} $X$]

\uline{Corollaire} 1.2 Soit $f:X\to Y$ une application
continue \deleted{propre}\inserted{fermée}, avec $Y$ loc compact,
\deleted{$F$ un} $y\in Y$, $X_0=f^{-1}(\{y\})$, $F$ un faisceau
abélien sur $X$, $F|X_0$ le faisceau induit
sur $X_0$. Alors \unsure{on a} un isom
canonique
\[
R^*f_{*}(F)_y\simeq H^*(X_{\unsure{0|y}},F|X_y)
\]
% archive sheet number: 247

\scanpage{249}\mspage{- \handwritten{6} -}
n$^{\mathrm{o}}$ \handwritten{293}

$\mathrm z$ soit l'ensemble des éléments $z\in\mathrm r$ tels que $zu\in\mathrm f$.
Posons $\mathrm r'=\mathrm o[t]+u\mathrm r$ ; on vérifie immédiatement que $\mathrm r'$
est un sous-anneau de $\mathrm r$. Si $y$ est un élément de $\mathrm o[t]$, une
condition nécessaire et suffisante pour que $y\,\mathrm r'\subset\mathrm o[t]$ est
que $yu\mathrm r\subset\mathrm o[t]$, donc que $yu\in\mathrm f$, ou encore que $y\in\mathrm z\cap\unsure{\mathrm e|\mathrm o}[t]={}
=\mathrm y$ ; $\mathrm y$ est donc bien le conducteur de $\mathrm r'$ par rapport à $\mathrm o[t]$.
Supposons pour un moment que $\omega(t)$ ne soit pas transcendant par
rapport à $\omega(\mathrm o)$ ; il y a alors un polynôme $\bar H\ne 0$ à coefficients
dans $\omega(\mathrm o)$ t\deleted{\unsure{a}}el que $\bar H(\omega(t))=0$. On en déduit tout de suite qu'il
y a un polynôme $H$ à coefficients dans $\mathrm o$ dont le coefficient dominant
$c$ n'est pas dans $\mathrm z$ tel que $H(t)\in\mathrm z$. On a alors $H(t)\in\mathrm y\unsure{'|\vphantom{'}}$ ;
il résulte alors du lemme \handwritten{3}, appliqué au couple $(\mathrm o,\mathrm r')$,
qu'il y a une puissance de $c$ qui appartient à $\mathrm y$, donc à $\mathrm z$ ; or
ceci est impossible puisque $\mathrm z$ est premier et $c\notin\mathrm z$.

\uline{Lemme} \handwritten{5} - La \uline{situation étant celle décrite dans l'énoncé du
lemme} \handwritten{3},- \uline{supposons de plus que} $\mathrm o$ \uline{soit noetherien}, \uline{Soient} $\mathrm q$ \uline{un
idéal premier maximal de} $\mathrm r$ \uline{et} $\mathrm m$ \uline{l'idéal premier} $\mathrm q\cap\mathrm o$ \uline{de} $\mathrm o$,
\uline{Si l'anneau,local de} $\mathrm q$ \uline{est distinct de celui de l'idéal} $\mathrm m$ \uline{de} $\mathrm o$,
\uline{il y a un idéal premier de} $\mathrm r$ \uline{contenu dans} $\mathrm q$, \uline{distinct de} $\mathrm q$
\uline{et contenant} $\mathrm m$.

Supposons d'abord que $\mathrm q$ contienne le conducteur $\mathrm f$ de $\mathrm r$
par rapport à $\mathrm o[t]$ ; nous allons voir qu'on est alors dans le
second cas. L'idéal \hspace{1.5em} contient un élément minimal $\mathrm r$ de l'ensemble
des idéaux premiers qui contiennent $\mathrm f$ ; $\mathrm z$ est donc un idéal
premier associé à $\mathrm f$. Soit $\omega$ l'homomorphisme canonique de $\mathrm r$
sur $\mathrm r/\mathrm z$ ; $\omega(t)$ est transcendant sur $\omega(\mathrm o)$ (lemme \handwritten{3}) et $\mathrm r/\mathrm z$
est entier sur $(\omega(\mathrm o))[\omega(t)]$. Le noyau de la restriction
% archive sheet number: 248

\scanpage{250}
Voici un résultat \deleted{\unsure{plus}} profond, fondamental
\deleted{\unsure{dans la th}} en géométrie analytique :

\uline{Lemme} 1.3. (Grauert). Soit $f:X\unsure{-|\to}Y$ un
morphisme propre d'espaces analytiques sur $\mathbb C$,
$F$ un faisceau coh. sur $X$. Alors

(i) les $R^if_{\unsure{*}}(\deleted{\unsure{E}}F)$ sont cohérents

(ii) Soit $y\in Y$, \deleted{\unsure{Alors le système}}\inserted{$Y_n=\unsure{\operatorname{Specan}|\operatorname{Spec\,an}}(\mathcal O_y/m_y^{n+1})$},
$X_n=X\times_Y Y_{\unsure{n|m}}$, $F_n=F\otimes_Y Y_n$, alors le système
projectif $(H^i(X_n,F_n))_n$ satisfait M.L.
et l'hom canonique
\[
R^if_*(F)_y{}^{\wedge}\longrightarrow\underset{n}{\underleftarrow{\lim}}H^i(X_n,F_n)
\]
est un isom .

\uline{Corollaire} 1.4. Supposons que
$y$ soit un point isolé de $\operatorname{supp} R^if_*(\unsure{E|F})$,
\deleted{Alors} i.e que $R^if_*(F)_y$ soit artinien.
Alors on a un isom
\[
R^if_*(F)_y\longrightarrow\underset{n}{\underleftarrow{\lim}}H^i(X_n,F_n)
\]
% archive sheet number: 249

\scanpage{251}\mspage{- \handwritten{2} -}
n$^{\mathrm{o}}$ \handwritten{293}

un corps des fractions de $\mathrm o/\mathrm p$ ; soit $R(\theta_1,..,\theta_n)$ une extension
transcendante pure de $R$ engendrée par $n$ éléments $\theta_1,..,\theta_n$ algébriquement
indépendants par rapport à $R$. L'homomorphisme canonique de
$\mathrm o$ sur $\mathrm o/\mathrm p$ se prolonge en un homomorphisme $\omega$ de $\mathrm o[t_1,..,t_n]$
dans $R[\theta_1,..,\theta_n]$ qui applique $t_i$ sur $\theta_i$ $(1\leq i\leq n)$ ; le noyau
de cet homomorphisme est $\mathrm{po}[t_1,..,t_n]$, comme il résulte tout
de suite du fait que $\theta_1,..,\theta_n$ sont algébriquement indépendants
sur $R$ ; par ailleurs, $\omega(\mathrm o[t_1,..,t_n])$ est un domaine d'intégrité
ce qui montre bien que $\mathrm{po}[t_1,..,t_n]$ est premier. De plus
$\omega(\mathrm o[t_1,..,T_n])$ n'est pas un corps. En effet, dans le cas contraire,
ce corps devrait contenir $R$, donc être identique à $R[\theta_1,..,\theta_n]$
or ce dernier anneau n'est pas un corps puisque $n>0$. O\deleted{\unsure{u|r}}n en conclut
que $\mathrm{po}[t_1,..,t_n]$ n'est pas un idéal premier maximal de
$\mathrm o[t_1,..,t_n]$. Par contre, $\mathrm q$ étant un idéal premier maximal de
$\mathrm r$, et $\mathrm r$ étant entier sur $\mathrm o[t_1,\ldots,t_n]$, $\mathrm q\cap\mathrm o[t_1,\ldots,t_n]$
est un idéal premier maximal de $\mathrm o[t_1,..,t_n]$ ; cet idéal contient
$\mathrm{po}[t_1,\ldots,t_n]$, et en est distinct en vertu de ce qu'on vient
de dire. Comme $\mathrm o[t_1,\ldots,t_n]$ est intégralement clos, on sait
qu'il y a un idéal premier $\mathrm n$ de $\mathrm r$ tel que $\mathrm n\subset\mathrm q$, $\mathrm n\cap\mathrm o[t_1,..,t_n]$
$=\mathrm{po}[t_1,\ldots,t_n]$ ; cet idéal est contenu dans $\mathrm q$, distinct de
$\mathrm q$ et contient $\mathrm p$.

\uline{Lemme} \handwritten{2} - \uline{Soient} $\mathrm o$ \uline{un domaine d'intégrité local}, $\mathrm m$ \uline{son idéal
premier maximal,} $t$ \uline{un élément d'un sur-corps du corps des fractions
de} $\mathrm o$ ; \uline{on suppose que} $\mathrm o$ \uline{est intégralement fermé dans} $\mathrm o[t]$,
\uline{que} $t\notin\mathrm o$, \uline{et que l'idéal engendré par} $\mathrm m$ \uline{dans} $\mathrm o[t]$ \uline{ne
contient pas} $1$. \uline{Si} $F$ \uline{est un polynôme à coefficients dans} $\mathrm o$ \uline{tel}
% archive sheet number: 250

\scanpage{252}
Considérons maintenant, pour tt $n$, la suite
exacte canonique
\[
0\unsure{-|\to}\mathbb Z_{X_n}\overset{\alpha}{\longrightarrow}\mathcal O_{X_n}\overset{\beta}{\longrightarrow}\mathcal O_{X_n}^{\unsure{\wedge|\times|*}}\longrightarrow 0
\]
d'où une \deleted{\unsure{homomorphe|homomorph}}\inserted{longue} suite

\diagram{\begin{tikzcd}[ampersand replacement=\&,column sep=tiny,row sep=small,cramped,cells={font=\scriptsize}]
{..}\arrow[r] \& H^i(X_n,\mathbb Z_{\deleted{\unsure{X_n}}})\arrow[r,"\alpha^i"]\arrow[d,phantom,"\rotatebox{90}{$\simeq$}"] \&
H^i(X_n,\mathcal O_{X_n})\arrow[r,"\beta^i"] \& H^i(X_n,\mathcal O_{X_n}^{\unsure{\times|*}}]\arrow[r,"\delta^i"] \&
H^{i+1}(X_n,\mathbb Z_{\deleted{\unsure{X_n}}})\arrow[r,"\alpha^{i+1}"]\arrow[d,phantom,"\rotatebox{90}{$=$}"] \&
H^{i+1}(X_n,\mathcal O_{X_n}\\
\& H^i(X_0,\mathbb Z_{\deleted{\unsure{X_0}}}). \& \& \& H^i(X_{\unsure{n|0}},\mathbb Z) \&
\end{tikzcd}}

Pour \deleted{\unsure{un}} $n$ variable on trouve un syst
\deleted{suite exacte} projectif de suites exactes.

Je dis qu'\uline{il reste exact} \deleted{\unsure{à la limite}}\inserted{en passant à}
$\underleftarrow{\lim}$. En effet, \inserted{en vertu de (EGA, $\unsure{0|O}_{\mathrm{III}}$)} \unsure{cela} résulte du fait
que les Images des hom $\alpha^i,\beta^i,\delta^i$ \deleted{\unsure{f}} satisfont
M.L. Pour celles de $\alpha^i$ \inserted{et $\beta^i$}, cela résulte
du fait que \deleted{\unsure{$H^i(X$}} les systèmes
projectifs $(H^i(X_n,\mathbb Z))_n$ et $(H^i(X_n,\mathcal O_{X_n}))_n$
eux-mêmes y satisfont [le premier
parce qu'il est ess. \unsure{constant|constt}, le deuxième
en vertu de 1.3. (ii)]. Pour \unsure{les}
$\operatorname{Im}\delta^i$, ou ce qui revient au $=$
pour \inserted{les} $\operatorname{Ker}(\alpha^{i+1})$, on note que les derniers
\deleted{\unsure{forment}} \deleted{\unsure{un}} un syst. projectif \unsure{essentiellement|essentiellt}
% archive sheet number: 251

\scanpage{253}
% archive sheet number: 252
\margin{\handwritten{Grothendieck}}% short separate horizontal impression below the inscription

\hfill n$^{\mathrm{o}}$ \handwritten{293}

\begin{center}
\uline{LE THEOREME PRINCIPAL DE Z A R I S K I}

I. - LEMMES PRELIMINAIRES
\end{center}

\uline{Lemme} 1 - \uline{Soient} $\mathrm r$ \uline{un domaine d'intégrité et} $\mathrm o$ \uline{un sous-anneau de} $\mathrm r$. \uline{On suppose qu'il existe des éléments} $t_1,..,t_n$ \uline{de} $\mathrm r$ $(n>0)$ \uline{qui sont algébriquement indépendants par rapport à} $\mathrm o$ \uline{et tels que} $\mathrm r$ \uline{soit entier sur} $\mathrm o[t_1,....,t_n]$. \uline{Soient} $\mathrm q$ \uline{un idéal premier maximal de} $\mathrm r$ \uline{et} $\mathrm p=\mathrm q\cap\mathrm o$. \uline{Il y a alors un idéal premier de} $\mathrm r$ \uline{contenu dans} $\mathrm q$, \uline{distinct de} $\mathrm q$ \uline{et contenant} $\mathrm p$.

Nous nous ramènerons d'abord au cas où $\mathrm o$ est intégralement clos. Soit $\mathrm o^*$ la clôture intégrale de $\mathrm o$, et soit $\mathrm r^*=\mathrm r[\mathrm o^*]$; les éléments $t_1,...,t_n$ sont algébriquement indépendants sur $\mathrm o^*$ et $\mathrm r^*$ est entier sur $\mathrm o^*[t_1,...,t_n]$. Soit $\mathrm q^*$ un idéal premier de $\mathrm r^*$ tel que $\mathrm q^*\cap\mathrm r=\mathrm q$; comme $\mathrm r^*$ est entier sur $\mathrm r$, $\mathrm q^*$ est un idéal premier maximal de $\mathrm r^*$. Soit $\mathrm p^*=\mathrm q^*\cap\mathrm o^*$. Supposons qu'il y ait un idéal premier $\mathrm n^*$ de $\mathrm r^*$ contenu dans $\mathrm q^*$, distinct de $\mathrm q^*$ et contenant $\mathrm p^*$. Soit $\mathrm n=\mathrm n^*\cap\mathrm r$; $\mathrm n$ est un idéal premier de $\mathrm r$; comme $\mathrm q^*\cap\mathrm r=\mathrm q$ on a $\mathrm n\subset\mathrm q$; il résulte du fait que $\mathrm n^*\ne\mathrm q^*$ et de ce que $\mathrm r^*$ est entier sur $\mathrm r$ que l'on a $\mathrm n\ne\mathrm q$. Comme $\mathrm n^*$ contient $\mathrm p^*=\mathrm q^*\cap\mathrm o^*$, il contient aussi $\mathrm q\cap\mathrm o=\mathrm p$, d'où $\mathrm p\subset\mathrm n$.

Supposons donc maintenant que $\mathrm o$ soit intégralement clos. On sait qu'il en est de même de $\mathrm o[t_1,..,t_n]$. L'idéal $\mathrm p\mathrm o[t_1,..,t_n]$ engendré par $\mathrm p$ dans $\mathrm o[t_1,..,t_n]$ est premier. Soit en effet $\mathrm R$

\scanpage{254}
% archive sheet number: 253
\unsure{constant|constante|constants}. En effet (\unsure{cf|réf} ?) $H^i(X_n,\mathbb Z)=\deleted{\unsure{H|N}}H^i(X_0,\mathbb Z)$ est un groupe de type fini, et son quotient par $\operatorname{Ker}\alpha^i$ est sans torsion [car \unsure{isom} à un sous-groupe de $H^i(X_n,\mathcal O_{X_n})$, qui est un vectoriel sur $\mathbb C$], \deleted{\unsure{c'est|d'}} il suffit donc de noter que pour $n$ grand, \unsure{le|les|la} \unsure{rang|rangs} desdits quotients \unsure{reste|vont} \unsure{constant|constants|constante}.

On obtient ainsi, compte tenu de 1.3 (ii), et du fait qu'un \unsure{ext.|cat.} d'un syst. M.L. par un syst. M.L. est ML

\uline{\deleted{\unsure{Corollaire}}\inserted{Proposition}} 1.5\quad Sous les conditions de 1.3,
(i) Le syst. projectif $(H^i(X_n,\mathcal O_{X_n}^{\times}))_n$ satisfait ML

(ii) On a une suite exacte
\[
H^i(X_0,\mathbb Z)\xrightarrow{\alpha_\infty^{(i)}}
\deleted{\unsure{H|R}}R^if_*(\mathcal O_X)_y^{\wedge}
\xrightarrow{\beta_\infty^{(i)}}\underset{\leftarrow}{\lim}H^i(X_n,\mathcal O_{X_n}^{\times})
\xrightarrow{\unsure{\partial|\delta}_\infty^{(i)}} H^{i+1}(X_0,\mathbb Z)
\xrightarrow{\alpha_\infty^{(i+1)}} R^{i+1}f_*(\mathcal O_X)_y^{\wedge}
\]
\deleted{\illegible{1}}

Notons qu'on a une suite exacte déduite de la suite exacte
\[
0\longrightarrow \deleted{\unsure{\mathcal O|0}}\mathbb Z_X\longrightarrow \uline{\mathcal O}_X\longrightarrow \mathcal O_X^{\times}\longrightarrow 0
\]
par application de $(R^if_*)_y$ et compte tenu de 1.2, les \inserted{fibres en $y$ des} termes $R^if_*(\mathbb Z_X)_{\unsure{y|x}}$ \deleted{se} s'identifient à $H^i(X_0,\mathbb Z)$. On obtient ainsi une

\scanpage{255}
% archive sheet number: 254
\deleted{Madame,\\

À la suite de votre lettre du 24.6.\\ ma femme et moi avons l'impression\\ que vous \unsure{conviendrez|conviendriez} en effet au genre de\\ travail q.}

\scanpage{256}
% archive sheet number: 255
diagramme de suites exactes
\diagram{\begin{tikzcd}[ampersand replacement=\&,column sep=small,row sep=small,cells={font=\scriptsize}]
{} \& H^i(X_0,\mathbb Z)\arrow[r]\arrow[d,"u^i"] \& R^if_*(\mathcal O_X)_y\arrow[r]\arrow[d,"v^i"] \& R^if_*(\mathcal O_X^\times)_y\arrow[r,"\partial^i"]\arrow[d,"w^i"] \& H^{i+1}(X_0,\mathbb Z)\arrow[r,"\alpha^{i+1}"]\arrow[d,"u^{i+1}"] \& R^{i+1}f_*(\mathcal O_X)_y^{\unsure{\times|\vphantom{x}}}\arrow[d,"v^{i+1}"]\\
(D)\ \biggl\{ \& H^i(X_0,\mathbb Z)\arrow[r,dash]\arrow[d] \& R^if_*(\mathcal O_X)_y^{\wedge}\arrow[r,dash]\arrow[d] \& \underset{\leftarrow}{\lim}H^i(X_n,\mathcal O_{X_n}^{\times})\arrow[r,dash,"\widehat{\partial}^{i}"]\arrow[d] \& H^{i+1}(X_0,\mathbb Z)\arrow[r,"\widehat\alpha^{i+1}"]\arrow[d] \& R^{i+1}f_*(\mathcal O_X)_y^{\wedge}\\
\deleted{\unsure{0|o}}\arrow[r,dash] \& \deleted{H^i(X_0,\mathbb Z)}\arrow[r] \& \deleted{H^i(X_n,\mathcal O_{X_n})_{\unsure{n\in\mathbb N|n\ge0}}}\arrow[r] \& \deleted{H^i(X_n,\mathcal O_{X_n}^*)_{\unsure{n\in\mathbb N|n\ge0}}}\arrow[r] \& \deleted{H^{i+1}(X_0,\mathbb Z)}\arrow[r] \& \deleted{\unsure{R|H}^{i+1}(X_n,\mathcal O_{X_n})_{\unsure{n\in\mathbb N|n}}}
\end{tikzcd}}

On en conclut, grâce au lemme des $\overline{5}$

\uline{Corollaire} 1.6.\quad a) $R^if_*(\mathcal O_X^\times)_y\to\underset{\leftarrow n}{\lim} H^i(X_n,\mathcal O_{X_n}^\times)$ est \uline{injectif}, et les deux mbs ont $=$ image ds $H^i(X_n,\mathcal O_{X_n}^*)$.

b) \deleted{\unsure{Supposons}} Pour que $R^if_*(\mathcal O_X)_y$ soit \unsure{artinien}, i.e. $y$ un pt isolé de $\operatorname{supp}R^if_*(\mathcal O_X)$. Alors \deleted{\unsure{Alors}} \inserted{$\underset{U\ni y}{\underrightarrow{\lim}}H^i(f^{-1}(U),\uline{\mathcal O}_U^*)={}$} $R^if_*(\uline{\mathcal O}_X^\times)\to\underset{\leftarrow}{\lim}H^i(X_n,\uline{\mathcal O}_{X_n}^\times)$ \deleted{soit} un \uline{isomorphisme}.
% margin note written sideways
\margin{\deleted{\unsure{Comme les systèmes projectifs des} \illegible{1}}}
% slanted
\margin{A \unsure{expliciter|contre-exples} : pour $n$ grand Néron-Severi $(X_0)$ \unsure{reste constant} c'est i.e. :
$\operatorname{Ker}(H^2(X_0,\mathbb Z)\longrightarrow R^2f_*(\mathcal O_X)_{\unsure{y|0}})$}

c) \deleted{\unsure{Proposition} \unsure{1.7}.} Le lemme des $\overline{5}$ montre également que $R^if_*(\uline{\mathcal O}_X^\times)\longrightarrow\underset{\leftarrow}{\lim}H^i(X_n,\mathcal O_{X_n}^\times)$ est toujours un monomorphisme.
\deleted{\unsure{Comme|Si} $\unsure{\beta|\delta|\xi}=0$, $\{\underset{\leftarrow}{\lim}H^i(X_n,\mathcal O_{X_n}^\times),i\}$, il existe un \unsure{syst. proj.} \illegible{2} $(\unsure{R|H}^i(X_n,\mathcal O_{X_n}^*))_n$ satisfait M.L. (A.R. \unsure{ceci} \illegible{1})}
\deleted{\unsure{Comme} $v^{i+1}$ est injectif, \unsure{on} $[\,\operatorname{Ker}\alpha^{i+1}=\operatorname{Ker}\widehat\alpha^{i+1}\,]$, donc $\operatorname{Im}\partial^i=\operatorname{Im}\widehat{\partial}^i$ donc pour}
% slanted
\margin{$g\in R^if_*(\mathcal O_X)_y$ ayant \unsure{$\overline{m}$ image|m. image|$=$ image} dans $H^{i+1}(X_0,\mathbb Z)$ on \unsure{conclut} que \unsure{$Dg$ et $g$} sont congrus mod $\operatorname{Im}H^i(X_{\unsure{U|n|0}},\mathcal O_{X_{\unsure{y|n}}})$,
\uuline{Or} $R^if_*(\mathcal O_{X_U})\to H^i(X_0,\mathcal O_{X_n})$ est surjectif, on peut corriger, O.K.}

Le cas le plus important pour nous est celui où \unsure{$\nu$|$i$|$v$}$=1$, qui donne le

\uline{Th.} 1.\deleted{\unsure{8}}7. Soit $f:X\to Y$ un morphisme propre d'espaces analytiques \uline{sur} $\mathbb C$, $y\in Y$,
\inserted{$Y_n=\operatorname{Specan}\mathcal O_y/\mathfrak m_y^{n+1}$,} $X_n=X\times_Y Y_n$, \deleted{\unsure{Alors}} posons
\[
\begin{aligned}
\operatorname{Pic}(X_y)&=\underset{U\ni y}{\underrightarrow{\lim}}\operatorname{Pic}(f^{-1}(U))\\
&=R^1f_*(\uline{\mathcal O}_X^\times)_y,\quad \deleted{\text{Alors considérons l'hom.}}\\
\operatorname{Pic}(\widehat X_y)&=\underset{\leftarrow n}{\lim}\operatorname{Pic}(X_n).
\end{aligned}
\]
